\section{Spin, color, flavor, and virtuality characters}

Let \(\jmath:\Gamma_{\rm TPK}\to\Gamma_{\rm TPK}\) be the involution that reverses
the orientation of each edge and permutes the fibers through
\(J_c:E_c\to E_{\jmath c}\). A spin character is a representation
\(\rho_{\rm rot}\) of the subgroup generated by \(\jmath\) on the sections;
its sign eigenspaces \(+1\) and \(-1\) are algebraic objects distinct
from a physical representation of the Lorentz group.

Let \(\mathcal C\) be the Abelian group of channel characters of the transport.
Two sections have the same \emph{arithmetic flavour} when their characters
\[
 \chi_s:\mathcal C\longrightarrow\CC^\times
\]
coincide. Identifying these classes with fermion flavors requires
a representation of the algebra of physical observables and is not introduced by
nomenclature.

Finally, let
\[
 \operatorname{Ext}_N(s)=
 \{s'\in\Gamma(\mathcal B_{\gamma'}):
 \gamma'|_{\{0,\ldots,N\}}=\gamma,\ s'|_{\{0,\ldots,N\}}=s\}.
\]
The section is stable from depth \(N\) onward when the restrictions
\(\operatorname{Ext}_{N+1}(s)\to\operatorname{Ext}_N(s)\) admit a compatible branch
at every subsequent level. It is called \emph{virtual up to
depth \(N\)} when that global branch does not yet exist, although
finite extensions survive. This definition belongs to the inverse system and does not
depend on a temporal metaphor.

On the affine plane \(\FF_3^2\), the linear functional
\[
 \kappa_3:\FF_3^2\longrightarrow\FF_3,
 \qquad \kappa_3(r,c)=r-c
\]
is surjective, and its kernel has dimension one. Consequently, its three
fibers are the affine classes of the kernel, and each contains exactly three
points. The same balance persists in every disjoint union of complete
copies of \(\FF_3^2\). This ternary partition is a self-contained algebraic
statement. Once the qutrit chart, the orientation, and the phase character have been fixed,
the \cref{chap:algebra-color-qutrit} constructs the representation and its algebra
\(\mathfrak{su}(3)\); the intertwiner with TPK transport and the physical
chromodynamic identification remain an interface.
